% Chinese paper overview and demonstration guide.
% Rebuild with: make -C paper briefing
\documentclass[UTF8,fontset=none,a4paper,11pt]{ctexart}
\usepackage[margin=20mm,headheight=15pt,headsep=7mm,footskip=11mm]{geometry}
\usepackage{fontspec}
\setmainfont{DejaVu Sans}
\setsansfont{DejaVu Sans}
\setmonofont{DejaVu Sans Mono}[Scale=0.88]
\setCJKmainfont{Noto Sans CJK SC}
\setCJKsansfont{Noto Sans CJK SC}
\setCJKmonofont{Noto Sans CJK SC}
\usepackage{amsmath,booktabs,array,tabularx,graphicx,xcolor}
\usepackage[most]{tcolorbox}
\usepackage{enumitem,setspace,fancyhdr,lastpage}
\definecolor{fpInk}{HTML}{172634}
\definecolor{fpBlue}{HTML}{1D62D0}
\definecolor{fpTeal}{HTML}{007E6A}
\usepackage[unicode,colorlinks=true,linkcolor=fpBlue,urlcolor=fpTeal]{hyperref}
\hypersetup{pdftitle={KVFold：文章结构与演示介绍},
 pdfsubject={中文汇报材料；依据 demo 第 8 版与已完成的控制台验证}}
\color{fpInk}
\setstretch{1.15}
\setlength{\parindent}{0pt}
\setlength{\parskip}{0.45em}
\setlength{\emergencystretch}{2em}
\raggedbottom
\setlist{leftmargin=1.6em,itemsep=0.3em,topsep=0.4em,parsep=0pt}
\ctexset{section={format=\Large\bfseries\color{fpBlue},name={},number=\arabic{section},
 beforeskip=0pt,afterskip=0.65em}}
\pagestyle{fancy}
\fancyhf{}
\fancyhead[L]{\small\color{fpInk!65}KVFold · Paper Overview and Demonstration}
\fancyhead[R]{\small\color{fpInk!65}2026-11-09}
\fancyfoot[C]{\small\color{fpInk!60}\thepage\ / \pageref*{LastPage}}
\renewcommand{\headrulewidth}{0.3pt}
\newcolumntype{L}[1]{>{\raggedright\arraybackslash}p{#1}}
\newcolumntype{Y}{>{\raggedright\arraybackslash}X}
\renewcommand{\arraystretch}{1.32}
\newtcolorbox{messagebox}{colback=fpBlue!4,colframe=fpBlue!35,
 boxrule=0.5pt,arc=1.5mm,left=3mm,right=3mm,top=2mm,bottom=2mm,
 before skip=0.5em,after skip=0.6em}
\newcommand{\source}[1]{{\footnotesize\color{fpInk!65}#1\par}}
\newcommand{\consoleimage}[2]{%
 \begin{center}
 \includegraphics[width=\linewidth,height=94mm,keepaspectratio]{#1}\par
 \vspace{1mm}{\small\color{fpInk!75}#2}
 \end{center}}
\newcommand{\nextpage}[1]{\clearpage\section{#1}}

\begin{document}

{\fontsize{25}{32}\selectfont\bfseries\color{fpBlue}KVFold\par}
{\Large\bfseries 文章结构与演示介绍\par}
\vspace{2mm}
{\small Expanding Cache Capacity for LLM Agents with Lossless Compression\par}

这篇 demo 文章围绕 Agent 反复使用长上下文时的 KV 容量问题展开：
通过无损压缩，让固定的 GPU 和外部缓存预算容纳更多 KV，
再用前缀复用、并发请求和历史会话三个场景展示它带来的变化。

\section*{文章怎样组织}

\begin{small}
\begin{tabularx}{\linewidth}{@{}L{4.25cm}Y@{}}
\toprule
Section & Main content \\
\midrule
I. Introduction &
从编程 Agent 继续任务的例子引出问题；交代前缀缓存、容量瓶颈和相关工作，
再概括 KVFold 的方案与主要结果。背景知识和研究现状整合在这一节。 \\
II. System Overview &
介绍两条压缩路径，以及 KV 的保存、恢复和 attention 计算。
Pool 在 GPU 中保存压缩记录；Tier 压缩外部缓存，保留原生 GPU 路径。 \\
III. Demonstration &
文章重点。通过三个控制台场景，依次展示前缀复用与恢复、
GPU 容量对并发的影响、外部容量对历史留存的影响。 \\
IV. Conclusion &
归纳固定预算下增加 KV 容量的效果，以及它对请求等待和历史复用的意义。 \\
\bottomrule
\end{tabularx}\par
\end{small}

\section*{演示怎样串起来}

从一个排查 timeout 的编程任务开始：先建立上下文，再追加回归测试任务；
随后增加同时运行的会话，最后回到之前的历史任务。
三个场景分别回答“哪些前缀能复用”“GPU 能容纳多少活跃上下文”
和“旧会话的 KV 还能否留在缓存中”。

\begin{messagebox}
场景一已通过真实模型和 LMCache 跑通，并保存了运行记录。
场景二、三展示论文中的历史容量实验；页面选择器切换已有实验配置。
\end{messagebox}

Pool 使用 gather + FA3 完成 prefill，并在 fused decode kernel 中读取压缩 KV。
Tier 保持原生 GPU KV 与 attention，只压缩 DRAM/NVMe 中的缓存记录。
这两个模式对应后面两类容量对比。

\source{材料：英文论文第 8 版；控制台验证记录与真实页面截图。}

\nextpage{场景一：继续任务，观察前缀复用}

观众依次使用 Build、Append、Edit early 和 Restore。
页面把操作与前缀复用条放在一起，并显示输入、生成结果、
cached/computed tokens 和 engine TTFT。

\consoleimage{images/console-prefix-full(8).png}
{Figure 1. Prefix workflow from the recorded functional run.}

\begin{small}
\begin{tabularx}{\linewidth}{@{}L{2.5cm}Yrr@{}}
\toprule
Action & Interaction & Cached & Computed \\
\midrule
Build & 提交初始 timeout 排查任务，建立上下文 & 0 & 534 \\
Append & 保留原上下文和回答，追加回归测试任务 & 512 & 144 \\
Edit early & 修改开头内容，观察后续前缀复用失效 & 0 & 534 \\
Restore & 清空 GPU 前缀缓存，再请求原任务 & 512 & 22 \\
\bottomrule
\end{tabularx}\par
\end{small}

Restore 已实际从 LMCache 的 CPU DRAM 取回 512 个 token。
模型 KV 检查器另取一个已编码的完整 block，经过 GPU 导出、主机搬运和
独立 GPU 区域导入后，全部层的 BF16 字节不一致数为 \textbf{0}。

\begin{messagebox}
这个场景让观众直接看到：追加任务可以复用前缀，靠前的修改会影响后续复用，
GPU 中的前缀清除后还可以从外部缓存恢复。
\end{messagebox}

\source{功能验证：Qwen3-4B-Instruct-2507，单 H800 PCIe，TP=1，BF16 KV；
4 GiB GPU KV、1 GiB LMCache DRAM。当前截图显示保存的真实运行；
本轮外部恢复使用 DRAM，未启用 NVMe。}

\nextpage{场景二：增加会话，观察 GPU 容量}

观众切换并发准入上限 $S=8$ 和 $S=12$，
比较同一 GPU KV 预算下的 Raw 与 Pool。页面同时显示可用 token 容量、
等待请求数、engine TTFT 和 goodput。

\consoleimage{images/console-gpu-full(8).png}
{Figure 2. Recorded GPU-capacity comparison at admission cap S=12.}

\begin{small}
\begin{tabularx}{\linewidth}{@{}Yrr@{}}
\toprule
Metric & Raw & Pool \\
\midrule
Usable KV capacity (M tokens) & 1.18 & 1.59 \\
Mean queue (requests), $S=12$ & 1.5 & 0.4 \\
Mean engine TTFT (s), $S=12$ & 5.16 & 1.94 \\
Goodput (requests/min), $S=12$ & 10.6 & 14.9 \\
\bottomrule
\end{tabularx}\par
\end{small}

演示先展示 $S=8$，再切到 $S=12$。Raw 的平均 TTFT 从 1.69 s 增至
5.16 s，Pool 从 1.34 s 增至 1.94 s。
同样 54 GiB/GPU KV 预算下，Pool 可容纳的 token 增加约 34.7\%；
在 $S=12$ 时，goodput 提高约 40.6\%。

\begin{messagebox}
这一页的重点是容量增加如何缓解请求等待。
四个配置的整体 cached-token share 均为 0.94，
对比突出的是 GPU 中能同时容纳多少 KV。
\end{messagebox}

\source{历史 AgentX 实验：Qwen3-30B-A3B-Instruct，TP=2，双 H800；
$P=64$，400 GiB DRAM、2 TiB NVMe。Engine TTFT 不含代理排队时间；
goodput 统计满足 TTFT $\leq$2 s、ITL $\leq$40 ms 的请求。}

\nextpage{场景三：回到旧会话，观察历史留存}

观众切换 640 GiB 和 2 TiB NVMe 预算，对比 Raw 与 Tier。
两者保持相同的原生 GPU KV 和 attention 路径，
观察压缩外部记录后能保留多少可复用历史。

\consoleimage{images/console-tier-full(8).png}
{Figure 3. Recorded external-cache comparison at 640 GiB NVMe.}

\begin{small}
\begin{tabularx}{\linewidth}{@{}L{2.7cm}Yrr@{}}
\toprule
NVMe budget & Metric & Raw & Tier \\
\midrule
640 GiB & Cached-token share & 0.55 & 0.91 \\
640 GiB & Mean engine TTFT (s) & 6.76 & 2.34 \\
640 GiB & Goodput (requests/min) & 4.6 & 13.5 \\
2 TiB & Cached-token share & 0.94 & 0.94 \\
2 TiB & Mean engine TTFT (s) & 1.69 & 1.85 \\
\bottomrule
\end{tabularx}\par
\end{small}

在 640 GiB 预算下，压缩让更多历史留在缓存中。
切到 2 TiB 后，两者都能保留工作集，Tier 的 TTFT 略高，
体现编解码开销。Cached-token share 统计 GPU 与外部缓存共同提供的前缀。

页面还用独立实验说明 DRAM 驻留：原始工作集为 261 GiB，压缩后为
188 GiB。在 220 GiB 预算下只有压缩数据能全部留在 DRAM；
切到 400 GiB 后两者都能放下。

\begin{messagebox}
先展示容量紧张时保留更多历史的收益，再展示原始数据已经放得下时的成本，
让观众理解压缩收益取决于具体的容量条件。
\end{messagebox}

\source{历史 AgentX 配置：$P=64$、$S=8$、400 GiB DRAM；
DRAM 工作集示例来自另一组 48-prompt 实验。
容量数据保留原报告的测量版本，早于最新码表与索引调整。}

\end{document}
